\cleardoublepage

\section{总结与展望}

\subsection{本文总结}

本文围绕多跳问答中的证据覆盖、系统成本和回答可信度，搭建并分析了一套可复现实验原型。与其把RAG简单理解为“检索后再生成”的直线流程，本文更关心检索什么时候需要补一次、补完之后证据有没有真的变得更可用，以及生成后的验证究竟是在提高可靠性，还是在额外制造新的错误。

从实现上看，本文主要完成了三件事：一是把文本、表格和图谱片段统一转写为可检索的序列化证据，并放入同一套检索与重排序框架中；二是引入基于重排序分数的Adaptive补检索门控，让系统在必要时再做一次查询，而不是对所有样本一视同仁地增加成本；三是在答案生成之后接入基于 CoVe 思想的启发式后验验证，以便观察拒答机制是否真的能改善系统可信度。相比单纯罗列模块，更重要的是本文把这些组件放进同一套实验体系里，逐项检查它们带来的收益和副作用。

实验结果说明，当前结论必须分层理解。纯文本基线仍然是最稳定的端到端配置；异构序列化在现阶段并没有稳定转化为收益，反而因为格式噪声拖低了F1；Adaptive 在最终答案分数上的帮助有限，但在检索覆盖率上确实提供了小而稳定的增益，尤其体现在 bridge 类问题上；CoVe 验证则暴露出最明显的系统边界，软阈值时约束不足，严格阈值时又导致大规模错误拒答。也就是说，本文得到的最重要结论不是“把所有模块叠起来以后系统更强”，而是当前这条轻量级RAG流水线在什么条件下还能工作、又是在什么地方开始明显失效。

\subsection{未来展望}

后续最值得继续推进的方向有三点。第一，应当把验证模块从单纯的“拦截器”改成真正能够回写检索的反馈节点。当前系统一旦验证失败，只会直接拒答，但很多失败样本其实暴露的是“证据还差半步”，而不是“根本没有答案”。如果能让验证失败反向触发补检索，系统才有机会从死板拒答走向自我修复。

第二，Adaptive 目前仍是启发式门控，虽然实现简单、解释也清楚，但它还没有显式学习“多检一次值不值”。未来可以把期望覆盖率增益、额外Token开销和时延损失一起建模，让系统学会在预算约束下做检索决策，而不是完全依赖人工设阈值。

第三，验证阶段需要更细粒度的语义判定。本文的实验已经说明，单靠近似字面匹配很难同时兼顾安全性和可用性。更合理的做法，是引入轻量级NLI或语义匹配模型，对“完全不支持”“部分支持”和“明显矛盾”作出区分。只有把这一步做好，拒答机制才可能真正成为系统可靠性的补充，而不是新的误差来源。
